PWC is not yet a standard metric in redistricting litigation, but it is gaining
traction as courts and redistricting commissions move toward representation-quality
criteria.

\subsection{Academic Foundation}

\citet{fryer2011measuring} provided the theoretical framework and showed that
PWC-minimising redistricting plans reduce the ``travel cost'' of political
participation for constituents.
Their formulation has been cited in academic redistricting literature and
is embedded in the ``community of interest'' operationalisation used by
several state redistricting commissions.

\subsection{VRA District Analysis}

PWC is particularly relevant for Voting Rights Act (VRA) analysis of
majority-minority districts.
A VRA-compliant district that concentrates minority voters near the district
centroid (low PWC) may be more legally defensible under the ``functional''
community-of-interest test than a district where minority voters are scattered
across a wide geographic area (high PWC).
The prime-factor algorithm's PWC advantage makes it particularly suitable
for VRA district analysis.

\subsection{State Redistricting Criteria}

Several state redistricting guidelines include ``community of interest'' as
a criterion~\citep{ncsl2021criteria}.
PWC provides a spatial operationalisation of this criterion: a district with
low PWC has its population clustered near a geographic centre, consistent with
the community-of-interest standard.
While no state has explicitly adopted PWC as a formal metric, the underlying
concept (residents clustered together) is widely used.

\subsection{Practitioner Recommendation}

Include PWC in the composite compactness profile (K.7) as the representational
dimension, alongside the five geometric metrics.
Frame PWC as a ``representation quality'' measure rather than a ``shape'' measure,
and note that prime-factor achieves the best PWC performance.
For VRA districts, PWC analysis of the minority-voter clustering should
accompany the standard geometric compactness analysis.
